\section{The \instname\ data reduction flow.}
\label{sec:reduction_flow}


The overall data flow of the \instname\ pipeline is displayed in Figure \ref{fig:reduction_cascade}.

 
\begin{figure}[h!]
\centering
\includegraphics[width=0.9\textwidth]{../instrument_tutorial/figures/reduction_cascade.jpg}
\caption{The data reduction cascade of the \instname\ workflow.}
\label{fig:reduction_cascade}
\end{figure}



The reduction cascade is organized in tasks, which represent
well-defined steps in the process. Tasks can be grouped inside
sub-workflows.  Each task runs a recipe; the detailed description of
the algorithms, input, outputs and recipe parameters used in each
recipe are available in the pipeline manual. Here, we present only the
description of most important features.

The \wkfname\  \edps\ workflow is designed to execute the tasks that
deliver the final reduced data cube for each dataset. It can be either
the product of a single exposure, or the combination of multiple
exposures. Only calibrations needed by the selected the scientific
exposures are processed.

It is possible to set \edps\ to perform the data reduction until a
certain step of the reduction chain (e.g. to reduce only standar
stars, or only flat fields).  This is done by specifying the desired
tasks in the field {\bf Select reduction target} of the {\bf Raw Data}
tab.

The reduction steps of the \wkfname\ workflow are listed below. Before
starting the reduction, the parameters of the recipes associated to
each task can be configured by pressing the button

\includegraphics[width=0.8cm]{../common/figures/configure_dataset.jpg}
close to each dataset configuration.
See for more info on the configuration editor \ref{sec:configuration_editor}



\subsection{Imaging workflow}
\subsubsection{Bias}

The task {\bf bias} reduces bias frames and generates a master bias
calibration. Typically, default parameters give satisfactory results.


\subsubsection{Sky flat}
The task {\bf bias} reduces sky flats and generates a master flat
field calibration. Typically, default parameters give satisfactory
results.  If the field does not contain bright stars, it might be
advisable to set to the 'stack\_method' to 'mean' fron the
configuration editor (Section \ref{sec:configuration_editor}).

\subsubsection{Photometric Standard}

The task {\bf photometric\_standard} runs the recipe {\tt
  efosc\_zeropoint} to process observations of standard stars and
compute the photometric zeropoint and exinction coefficient. The
recipe runs sextractor to identify objects and perform aperture
photometry. Typically, default recipe parameters deliver good
results. Standard stars are used to determine the zeropoint only for
filters “U\#640” (using catalogue from Landolt 1992, 2007), and filters
“B\#639”, “V\#641”, “R\#642”, “i\#705” (using catalogs from Stetson 2000).

This extinction coefficient may not be the correct one for a given
night. In such a case the zeropoint will reflect the mismatch of the
extinction. If the airmass differs between the observations of the
standard star and that of the science the mismatch will not be fully
corrected by the zeropoint. Please refer to the EFOSC2 pipeline user
manual (Izzo et al. 2016: Sections 9 and 10) for the details of this
recipe and the extraction algorithms employed.

Note that the results of this task are not used to calibrate the
scientific exposure.


\subsubsection{Science image}

The task {\bf science\_imaging} runs the recipe {\bf
  efosc\_img\_science} to reduce the raw image, and compiles a
catalogue of extracted sources using sextractor.

If zeropoint and exinction coefficient determined from the standard
star have to be applied, they have to be applied outside the workflow, as the feature is not yet supported by the pipeline.


\subsection{Spectroscopic workflow}


\subsubsection{Bias}

The task {\bf bias} reduces bias frames and generate a master bias
calibration. Typically, default parameters give satisfactory results.


\subsubsection{Spectroscopic calibrations}

The task {\bf Spectroscopic\ calibrations} runs the recipe {\tt
  efosc\_calib} to process arc lamps and flat fiels to determine the
wavelength calibration, spectral distortion, and the flat fielding of
spectroscopic calibrations. All spectroscopic modes are supported.
Typically, default parameters give satisfactory results.

\subsubsection{Spectroscopic standard}

The task {\bf spectroscopic\_standard} runs the recipe {\tt
  efosc\_science} on the standar stars and produces a response curve,
which will subsequently be used to flux-calibrate the science
observation. Typically, default recipe parameters deliver good
results.

\subsubsection{Science spectrum}

The task {\bf Science spectrum} reduces scinetific spectroscopic
observations for all spectroscopic modes. The \instname\ workflow will
flux-calibrate the science observation using the instrument response
curve derived from the standard star observation if it exists in the
current DataSet. If no standard star observation exists in the current
DataSet, then the science observation will not be flux-calibrated.
